
\documentclass[11pt]{article}
\usepackage[margin=1in]{geometry}
\usepackage[T1]{fontenc}
\usepackage[utf8]{inputenc}
\usepackage{amsmath,amssymb,amsthm}
\usepackage{booktabs}
\usepackage{array}
\usepackage{enumitem}
\usepackage{xcolor}
\usepackage{hyperref}
\usepackage{url}
\hypersetup{colorlinks=true,linkcolor=blue!50!black,citecolor=blue!50!black,urlcolor=blue!50!black}
\setlist[itemize]{leftmargin=1.4em,itemsep=0.25em,topsep=0.25em}
\setlist[enumerate]{leftmargin=1.6em,itemsep=0.25em,topsep=0.25em}
\newcommand{\WMPLN}{WM--PLN}
\newcommand{\PLN}{PLN}
\newcommand{\KS}{K\&S}
\newcommand{\codepath}[1]{\path{#1}}
\newcommand{\leanref}[1]{\path{#1}}
\newcommand{\Evidence}{\ensuremath{\mathcal{E}}}
\newcommand{\World}{\ensuremath{\mathcal{W}}}
\newcommand{\Query}{\ensuremath{\mathcal{Q}}}
\newcommand{\R}{\ensuremath{\mathbb{R}}}
\newcommand{\draftnote}{\noindent\textit{Draft extended abstract for author revision.}}
\sloppy

% Draft inspiration only. Author should rewrite, verify claims, and choose final voice.
\title{Updating Causal World Models with Evidence-Carrying Revision}
\author{Zarathustra Amareis Goertzel\\\small Mettapedia Project}
\date{Draft for JoeFest extended abstract}
\begin{document}
\maketitle
\draftnote

\begin{abstract}
Recent work by Halpern and collaborators studies agents who are uncertain about the true causal model, choose actions according to their current causal beliefs, and then update those beliefs using feedback generated by their own choices. This extended abstract sketches a WM--PLN formalization of that situation. A world-model state carries evidence over candidate structural causal models, interventions produce observations, revision updates the evidence state, and queries extract posterior support, confidence, and explanatory summaries. The framework naturally distinguishes ordinary Bayesian updating, inexplicable observations that require model-support expansion, and self-confirming steady states where an agent's actions generate evidence compatible with false causal beliefs. The proposed experiment compares greedy action selection with evidence-diversifying exploration. The goal is to bring Halpern's semantics of causality, uncertainty, and decision into an evidence-carrying PLN architecture.
\end{abstract}

\section{Halpern's causal-belief problem}

Halpern and Pearl's structural-model account gave a precise semantics for causal explanation and actual causality \cite{HalpernPearl2005,Pearl2000}. More recent work with Piermont and Vier\o{} asks a complementary question: how can a decision maker's subjective causal judgments be recovered from preferences over interventions, and how should causal beliefs update when an agent acts and observes the consequences? \cite{HalpernPiermont2024,HalpernPiermontViero2026}.

This is a natural target for WM--PLN. PLN traditionally emphasizes uncertain inference over propositions and links. The world-model calculus generalizes the state being revised: the uncertain object can be not only a fact, but a causal model, a source-reliability model, or a policy for updating in the face of anomaly.

\section{A finite causal world-model interface}

Let $\mathcal{M}$ be a finite set of candidate structural causal models, $A$ a set of interventions, $O$ a set of possible observations, and $U$ a utility function on outcomes. A WM--PLN causal state can be represented abstractly as an evidence-weighted distribution over models:
\[
  W : \mathcal{M} \to \Evidence.
\]
Queries include:
\begin{itemize}
\item model-support queries: ``how much evidence supports model $M_i$?'';
\item intervention queries: ``what is the predicted outcome distribution under $do(a)$?'';
\item explanation queries: ``which causal model families explain observation $o$ after action $a$?'';
\item policy queries: ``which intervention maximizes expected utility under the current evidence state?''
\end{itemize}

Revision is indexed by action and observation:
\[
  \mathrm{revise}(W,a,o) = W'.
\]
In the ordinary explainable case, this is Bayesian or conjugate evidence update over candidate models. In the anomalous case, where $o$ has zero or negligible likelihood under the current support, the system must either expand model support, introduce an unknown-mechanism bucket, or invoke a trust-gated anomaly policy.

\section{Three tribute experiments}

\paragraph{1. Unintended consequences.}
The agent believes an intervention will reduce an undesirable outcome, but the intervention changes incentives or mechanisms so that the outcome worsens. The classic ``cobra bounty'' pattern is a vivid test case: an action selected under a plausible causal model generates evidence that cannot be rationalized by that model. WM--PLN should show the old model losing support or the hypothesis space expanding.

\paragraph{2. Self-confirming delusion.}
The agent holds a false causal belief but repeatedly chooses an action whose feedback is compatible with that belief. Naive revision can converge to a stable but false evidence state. This is a world-model analogue of a steady state: the agent's optimal action, given its beliefs, generates observations rationalized by those same beliefs.

\paragraph{3. Exploration repair.}
Add a control policy that values expected model discrimination or evidence diversity. The policy may sacrifice immediate expected utility in order to sample actions that distinguish candidate causal models. The experiment then compares greedy expected-utility choice with evidence-diversifying choice.

\section{Why WM--PLN is a useful lens}

The evidence-carrying view adds three ingredients to ordinary causal-model uncertainty.

First, it separates posterior support from confidence. Two causal models may have similar posterior probabilities but very different effective evidence counts. This matters for whether the system should exploit, explore, or abstain.

Second, it supports multiple update regimes. Additive revision handles independent evidence; trust-gated revision handles source conflict; overlap-aware revision prevents double-counting; model-support expansion handles genuinely surprising observations.

Third, it makes local PLN chaining a subroutine rather than the whole semantics. Within a fixed causal model, PLN-style chaining can approximate or compile local causal reasoning. When the model itself is uncertain or misspecified, the evolving world-model state is the proper object of inference.

\section{Expected contribution}

This version of the extended abstract is more experimental than the semantic-state and chaining versions. Its value is that it directly celebrates Halpern's late-career synthesis of causality, uncertainty, action, and belief update. The resulting WM--PLN experiment would show not merely how uncertain conclusions are propagated, but how the agent's model of the causal world is revised by the consequences of its own actions.

\begin{thebibliography}{99}

\bibitem{HalpernPiermont2024}
J.~Y. Halpern and E.~Piermont.
\newblock A representation theorem for causal decision making.
\newblock In \emph{Proceedings of the 21st International Conference on Principles of Knowledge Representation and Reasoning (KR 2024)}, pages 410--419, 2024.
\newblock DOI: 10.24963/kr.2024/39.

\bibitem{HalpernPiermontViero2026}
J.~Y. Halpern, E.~Piermont, and M.-L. Vier\o{}.
\newblock Unintended consequences: Updating causal models.
\newblock arXiv:2603.09387, 2026.

\bibitem{Halpern2003}
J.~Y. Halpern.
\newblock \emph{Reasoning about Uncertainty}.
\newblock MIT Press, 2003.

\bibitem{HalpernPearl2005}
J.~Y. Halpern and J.~Pearl.
\newblock Causes and explanations: A structural-model approach. Parts I and II.
\newblock \emph{British Journal for the Philosophy of Science}, 56(4):843--887 and 56(4):889--911, 2005.

\bibitem{Pearl2000}
J.~Pearl.
\newblock \emph{Causality: Models, Reasoning, and Inference}.
\newblock Cambridge University Press, 2000.

\bibitem{GoertzelPLN2009}
B.~Goertzel, M.~Ikl\'{e}, I.~F. Goertzel, and A.~Heljakka.
\newblock \emph{Probabilistic Logic Networks: A Comprehensive Framework for Uncertain Inference}.
\newblock Springer, 2009.

\end{thebibliography}
\end{document}
